\chapter{终端定位算法相关理论与技术基础}

\section{传统无线电定位技术理论}

在详细探讨本文提出的深度学习与强化学习定位框架之前，有必要对经典的无线电定位技术进行理论回顾。传统的无线电参数估计定位技术主要基于接收信号的空间几何特征（如角度、时间、电平强度）来实现目标位置的解算。其主流方案可归纳为以下几类：

\subsection{基于到达角（AOA）的定位技术}

到达角（Angle of Arrival, AOA）定位算法主要依靠多天线阵列测量目标发射信号到达各个接收节点的空间入射方向。在二维平面定位中，至少需要两个相距一定距离的接收机节点。各个接收机通过阵列信号处理算法（如 MUSIC 或 ESPRIT 估计算法）测算出信号到达的方位角，并以自身为原点沿该角度作射线，多条射线的几何交点即为目标的空间坐标估计值。

AOA 技术的优势在于无需接收端与发射端之间进行高精度的时钟同步，同时测量所需的辅助节点数量较少（二维定位仅需两台设备）。然而，该方法硬性要求接收系统必须配备复杂的天线阵列以实现测向。更致命的是，AOA 定位精度对角度的微小测量误差极其敏感，在存在严重多径散射的城市环境中，直射波被遮挡，接收端捕获的多为经过多次反射的畸变波阵面，导致测角结果与真实物理方向发生严重偏离。

\subsection{基于到达时间（TOA）的定位技术}

到达时间（Time of Arrival, TOA）定位是一种基于绝对距离测量的几何定位方法。电磁波以光速 $c$ 传播，通过精确记录无线信号从发射机到达接收机所经历的绝对时间差 $\Delta t$，即可计算出两者之间的直线物理距离 $d = c \cdot \Delta t$。在二维平面内，至少需要三个位置已知的接收节点测得至目标的距离，随后以各节点为圆心、测量距离为半径作圆，三个圆的唯一交点即确定了目标的位置。

TOA 定位理论上能够提供极高的定位精度，但其在工程实现上面临着苛刻的同步挑战。由于光速极大，纳秒级的时间计量误差就会导致数米的测距偏差，因此系统要求所有接收节点以及作为被测目标的非法终端必须保持严密的纳秒级全局时钟同步。在非合作式的非法终端探测场景中，获取目标的同步时钟信息几乎是不可能的，这使得 TOA 原理难以直接应用。

\subsection{基于到达时间差（TDOA）的定位技术}

为克服 TOA 对发射端同步的苛刻要求，到达时间差（Time Difference of Arrival, TDOA）算法被广泛提出。TDOA 基本原理不再测量绝对飞行时间，而是测量目标信号到达不同接收节点（如节点 $i$ 与参考节点 $1$）的时间差 $\Delta t_{i,1}$。将时间差转化为距离差 $\Delta d_{i,1} = c \cdot \Delta t_{i,1}$。

根据解析几何原理，到两个定点距离之差为常数的点轨迹是一条双曲线。因此，通过一对节点可构造出一条以它们为焦点的双曲线。对于二维定位，至少需要布置三个相互独立的接收节点来联立求解两条双曲线的交点，以此逼近目标真实位置。TDOA 定位无需目标终端配合时钟同步，仅要求接收网络内部维持时间基准的一致性，因此广泛应用于合作与非合作无源定位系统中。但在城市遮挡环境下，非视距（NLOS）传播会给不同链路带来不同程度的额外时延，造成双曲线方程产生巨大偏差甚至无法相交。

\subsection{传统技术的局限性与突破思路}

综上所述，传统的 AOA、TOA 与 TDOA 等纯几何定位技术在视距（LOS）开阔空域表现优异，但在城市复杂电磁环境中均面临共同的技术瓶颈：其理论成立的根本前提是认为信号沿直线视距到达，一旦发生多径反射与 NLOS 阻挡，基于几何参数的直接解算便会发生系统性失效。这也正是本论文后续要探讨通过\textbf{射线追踪的高保真信道建模（第三章）}与\textbf{不依赖几何解算的深度神经网络（第四、五章）}来替代传统测向测距算法的理论出发点。

\section{卫星-无人机-终端立体定位架构}

\subsection{LEO卫星信号特性与定位潜力}

天地一体化网络是新一代移动通信技术的重要发展维度。近年来，以 Starlink 为代表的低轨（Low Earth Orbit, LEO）卫星通信星座迎来了爆发式增长。相比于传统的中高轨卫星，LEO 卫星运行在距离地面 $300 \sim 1500$ km 的轨道上，具有低传输时延、低路径损耗以及全球无缝覆盖等显著优势。

在定位应用层面，LEO 卫星宽带通信系统展现出了巨大的潜力。以 Starlink 系统为例，其下行链路工作在 Ku 波段（$10.7 \sim 12.7$ GHz），采用了正交频分复用（Orthogonal Frequency Division Multiplexing, OFDM）波形。在物理层协议规范中，Starlink 的无线帧（Frame）具有典型的时间-频率分布特征。一个完整的无线帧包含主同步信号（Primary Synchronization Signal, PSS）、辅同步信号（Secondary Synchronization Signal, SSS）以及公共参考信号（Common Synchronization Signal, CSS）。

其中，PSS 序列在时域上由 8 个重复的子序列构成，且首个子序列进行了反转处理。这种独特的重复结构设计，为接收端提供了极高的自相关处理增益，使其非常适合用于低信噪比（Signal-to-Noise Ratio, SNR）环境下的信号检测与时间同步。高带宽和特殊的帧结构意味着接收信号具有很高的时间分辨率，这为本课题提取信号级的指纹特征（如 I/Q 序列、信道脉冲响应CIR）提供了扎实的物理层理论支持。

\subsection{无人机辅助定位平台架构}

在城市环境中，非法卫星终端往往隐藏在建筑群中。传统的地面无线电监测手段存在两个致命缺陷：一是地面监测站的视野严重受限，极易被周围建筑物遮挡，导致完全处于非视距（Non-Line-of-Sight, NLOS）环境中；二是地面移动监测车的机动性较差，难以快速穿越拥堵的城市街道进行近距离定位。

为克服上述问题，本研究引入低成本无人机（Unmanned Aerial Vehicle, UAV）作为机载接收平台，构建了“卫星-无人机-终端”立体定位架构。无人机平台具有以下核心优势：
\begin{itemize}
    \item \textbf{高空视距优势：} 无人机可升至数十至数百米的高空，能够有效越过低层和中层建筑物的遮挡，在宏观层面上大幅改善监测设备与目标终端之间的视距（Line-of-Sight, LOS）条件。
    \item \textbf{三维高机动性：} 无人机不受地面交通地形限制，可以敏捷地调整飞行高度、悬停位置与飞行轨迹，从而从多个空间角度对目标终端进行信号探测。
    \item \textbf{低成本与易部署：} 随着消费级与工业级无人机技术的成熟，通过挂载轻量化的软件无线电（Software Defined Radio, SDR）接收机，即可快速形成一个兼顾经济性与灵活性的定位节点。
\end{itemize}
该立体架构为后续复杂城市环境下的多模态融合定位和多目标检测提供了理想的硬件承载平台。

\section{复杂城市环境无线电传播机理}

\subsection{城市多径效应与非视距传播机制}

无线电定位的精度从根本上受限于电磁波在空间中的传播环境。在密集的城市环境中，建筑物、植被、车辆等障碍物阻断了发射端与接收端之间的直射链路。当卫星发出的高频微波信号射向地面终端，或终端泄露的射频信号辐射向空中时，会与环境中的粗糙表面或尖锐边缘发生复杂的电磁相互作用。

主要的传播机制包括：
\begin{itemize}
    \item \textbf{反射（Reflection）：} 当电磁波照射到尺寸远大于其波长的平滑建筑表面（如玻璃幕墙、混凝土地面）时发生。
    \item \textbf{绕射（Diffraction）：} 当电磁波在传播路径上遇到锐利的物体边缘或尖角（如建筑物屋顶边缘）时，波面发生弯曲并扩散到障碍物的阴影区域。
    \item \textbf{散射（Scattering）：} 当电磁波遇到粗糙表面或尺寸小于波长的物体（如树叶、路灯杆）时，能量被向多个方向重新分配。
\end{itemize}

这些相互作用导致同一时间发射的信号，经过不同长度的路径，在不同时间到达接收天线，即所谓的\textbf{多径效应（Multipath Effect）}。在多径传播和 NLOS 条件下，接收到的信号是直射波与多条延迟、衰减、相位旋转的反射波和绕射波的叠加。这种严重的畸变使得传统的基于几何参数提取的定位算法（如到达角 AoA 和到达时间 ToA）产生极大的测量误差，进而导致定位解算完全失效。

\subsection{基于射线追踪的确定性信道建模理论}

为了深入研究复杂城市场景下的定位技术，首先必须对电磁信道进行高精度的建模。传统的信道模型主要是基于概率统计经验模型（如 Rayleigh 衰落模型或 3GPP 城市微小区模型）。尽管计算简便，但统计模型只能反映信号衰落的宏观平均分布特性，无法刻画某一特定三维坐标点上，信号具体的到达路径、角度和时延分布。对于以精确“位置特征”为依据的射频指纹定位而言，统计模型是完全失效的。

为了解决这一问题，本研究采用基于\textbf{射线追踪（Ray Tracing, RT）}的确定性信道建模理论。射线追踪是一种基于计算电磁学的高频渐近方法，其理论基础是\textbf{几何光学（Geometrical Optics, GO）}与\textbf{一致性绕射理论（Uniform Theory of Diffraction, UTD）}。
\begin{itemize}
    \item \textbf{几何光学（GO）：} 将电磁波抽象为沿直线传播的“射线（Ray）”，遵循斯涅尔定律计算射线在不同介质表面的入射、反射和折射系数。
    \item \textbf{一致性绕射理论（UTD）：} 用于弥补 GO 在处理阴影边界处的奇异性缺陷，通过引入绕射系数来精确计算电磁波在楔形或刃形边缘处的能量绕射。
\end{itemize}

利用射线追踪理论，并结合数字孪生（Digital Twin）技术导入真实的城市建筑物三维轮廓与材料电磁参数，可以在计算机中精确穷举出信号从源点到接收点的每一次反弹路径。该技术能够精确获取包含了多径时延、功率、相位、到达方位角和俯仰角等高分辨率特性的信道脉冲响应（Channel Impulse Response, CIR）。这种极其精确的空间几何-电磁映射关系，为第三章构建海量深度定位特征库和第五章生成高保真无线电热力图提供了坚实的理论依据与数据基石。

\section{深度学习在无线电定位中的应用基础}

\subsection{射频信号的深度特征表达}

在严重的城市多径干扰下，接收到的射频基带信号（如 I/Q 序列）发生了严重的相位偏转和多径混叠。传统的参数估计方法（如 MUSIC 算法）需要构建精确的方向向量或信号包络模型，在低信噪比和高度时变的 NLOS 环境中难以获得可靠特征。

深度学习中的卷积神经网络（Convolutional Neural Network, CNN）以其强大的非线性映射和局部特征感知能力，为绕开参数估计提供了一条全新捷径。一维 CNN 可以直接作用于时域的 I/Q 序列，通过多层卷积核提取出信号包络和相位中的隐式规律。进一步地，深度残差网络（Residual Network, ResNet）通过引入“残差块（Residual Block）”结构和跳跃连接（Skip Connection）机制，使得网络可以训练得更深且不会出现梯度消失问题。在无线电定位中，深层 ResNet 能够从海量受干扰严重的基带信号中，过滤掉瞬时的高斯自噪声，提取出与具体地理位置强相关的高维、抽象的环境指纹（Fingerprint）特征，大幅提升了极端传播环境下的定位鲁棒性。

\subsection{信号空间到图像空间的目标检测转化}

针对多目标并发的定位需求，传统的基于阵列天线的串行分离算法计算复杂度呈指数级增长，无法满足无人机平台的实时性要求。近年来，目标检测算法在计算机视觉领域取得了突破，其中 YOLO（You Only Look Once）系列以其单阶段检测带来的极致速度成为实时检测的标杆。

将信号处理转化为图像处理的核心理论基础在于\textbf{空间平滑与特征映射}。通过对机载接收机截获的空间离散信号强度（RSSI）应用高斯平滑滤波器进行空间插值，可以将零散的无线电采样数据映射重构为一张平滑的二维归一化“无线电热力图（Radio Map）”。在这一图像空间中，真实的非法终端呈现为符合高斯衰减规律的高亮能量聚集区（光斑），而多径干扰造成的假信号则表现为弥散、不规则的纹理。基于深度学习构建的 YOLO 模型（特别是摒弃特定先验锚框的 Anchor-Free 检测范式），通过深层卷积网络强大的纹理特征识别能力，能够一次性在热力图上回归出所有光斑的中心坐标。这种跨模态的技术迁移，从理论上将多目标通信定位的时间复杂度从 $\mathcal{O}(N^3)$ 降解为与目标数量无关的 $\mathcal{O}(1)$。

\section{强化学习在定位多源融合中的应用理论}

\subsection{多源融合定位的马尔可夫决策过程建模}

在复杂的城市飞行轨迹中，没有任何一种单一的测量模态能够保持绝对可靠：RSSI 极易受阴影衰落影响，AoA 极易受多径散射影响，而惯性导航（PDR）虽然短期平滑但存在长期累积误差。为了实现自适应的多模融合定位，传统的扩展卡尔曼滤波（EKF）需要预设固定的过程噪声和测量噪声协方差矩阵，这在时变的非平稳环境中容易导致滤波器发散。

为了突破静态滤波的限制，本课题引入了强化学习理论，将多模态权重的动态调度问题形式化为一个\textbf{马尔可夫决策过程（Markov Decision Process, MDP）}。MDP 由一个元组 $(S, A, P, R, \gamma)$ 组成。在融合定位的语境下，系统的“状态空间 $S$”不仅包括当前的定位解算结果，更应该包含当前信道的物理感知参数（如瞬时信噪比或时延扩展）；“动作空间 $A$”定义为对各模态融合权重的增减调节；“转移概率 $P$”由信道状态的时变性与无人机的物理运动学方程共同决定；而“奖励函数 $R$”则由定位精度误差和轨迹平滑度反馈构成。

\subsection{深度Q网络（DQN）自适应决策机制}

Q-Learning 是强化学习中经典的无模型算法，它通过更新一个作用于离散状态和动作的 Q 矩阵来寻找最优策略。但在高维且连续的通信信道特征面前，庞大的状态维度会导致“维数灾难”。

深度 Q 网络（Deep Q-Network, DQN）结合了深度神经网络对高维状态的表征能力和 Q-Learning 的决策机制。通过利用神经网络近似动作价值函数 $Q(s,a)$，DQN 智能体能够在不需要建立任何先验噪声模型的情况下，通过在数字孪生仿真环境中的大规模“试错（Trial and Error）”，自主学习识别信道质量规律。例如，智能体可以学习到：当状态空间中的多径时延扩展增大时，意味着当前处于严重的 NLOS 区域，应当降低射频模态（AoA/RSSI）的权重，而增加惯导 PDR 的权重。这种由物理环境特征驱动的自适应决策理论，构成了第四章 PA-DQN 融合框架的核心基础。

\section{空间分集增强低信噪比定位的理论基础}

\subsection{MIMO系统空间自由度分析}

多输入多输出（Multiple-Input Multiple-Output, MIMO）技术是现代通信系统的核心支柱。对于只有一根发射天线的非法终端被动探测定位系统，机载接收端的多天线阵列构成了单输入多输出（SIMO）系统，依然保留了 MIMO 技术中接收端的空间特性。

在数学机理上，一个具有 $N$ 根接收天线的阵列提供了 $N$ 个独立的\textbf{空间自由度（Spatial Degrees of Freedom）}。当接收不同到达角的射频信号时，各天线单元之间存在与物理空间排布相关的相位差。这些相位信息不仅被经典的波达方向（DoA）算法所利用，同时也为深度指纹网络提供了蕴含空间几何结构的本质差异特征，从而使得深度网络提取出的指纹信息具有更清晰的可分边界。

\subsection{多天线接收抗多径衰落机理}

不仅如此，由于空间中存在多个独立或部分相关的传播路径，信号在同一时间到达不同的天线单元时，其所经历的衰落情况通常是不同的。这就为接收端提供了\textbf{空间分集增益（Spatial Diversity Gain）}。

其理论机制在于，如果各天线单元的间距配置合理（通常大于半个波长），各天线接收到的信号深衰落谷底在空间上就不太可能完全重合。通过最大比合并（Maximum Ratio Combining, MRC）等合并策略，或者通过深度神经网络在底层特征映射中的自动加权聚合，可以大幅降低整个系统遭遇严重衰落的概率。正是这种由于物理空间分集所带来的抗多径衰落机理，从理论上保证了即使在信噪比极低（如 $-15$ dB）近乎被背景噪声淹没的极限工况下，仍能从多天线阵列中提取出支撑准确定位的有效射频特征。

\section{本章小结}

本章详细阐述了非法卫星终端定位技术研究中涉及的关键理论体系。首先解析了以 Starlink 为代表的 LEO 卫星通信体系，并分析了利用具备高空视野的挂载无人机构建“卫星-无人机-终端”立体监测架构的可行性与优势。其次，针对城市复杂遮挡问题，深入剖析了由反射、绕射等机制引起的多径传播与 NLOS 衰落痛点，并证明了基于几何光学与 UTD 的射线追踪理论是高保真数字孪生信道建模的最佳解决方案。随后，分别从深度学习的非线性抽象表征能力、目标检测在信号处理域的跨模态映射、强化学习 MDP 框架中 DQN 自主学习机制的动态权重调度优势，以及 MIMO 系统空间分集抗深衰落机理方面，构建了完整的相关理论基础。这些理论为后续构建第三章的深度特征提取模型、第四章的物理感知自适应融合网络，以及第五章的大并发目标成像定位体系提供了充分且严谨的地基。
